\begin{frame}{20 batches at each fan-out compare restore with a warm cached template.\ev{\tagP}}
\begin{tabularx}{\linewidth}{L{2.1cm}YY}\toprule
 & Warm template & Restore \\\midrule
Preparation & Dependencies baked in; pinned source & Capture after setup and one warm test \\
Fan-out & \multicolumn{2}{l}{$N\in\{3,6,12\}$; 20 interleaved batches per arm and $N$} \\
Endpoint & \multicolumn{2}{l}{Batch request $\rightarrow$ last child's first test result} \\\bottomrule
\end{tabularx}
\vspace{.25cm}
\begin{ticks}
\item No model calls; cold start is a reference.
\item Establish declared restore fidelity before timing.
\item Report capture cost, failed restores and batch latency separately.
\end{ticks}
\claim{A timing win supports this restore path on this workload.}
\sources{Experiment design: H1 in PREREGISTRATION.md; matched state and the last-child endpoint.}
\qadetail{Current protocol: gain is median cached makespan minus median restore makespan, with 10,000 within-arm bootstrap resamples. Supports H1 if the 95\% lower confidence bound exceeds 10 seconds at every N; rejects if the Bonferroni 98.3\% upper bound is below 10 seconds at any N. Gate 0 restores a RAM-only nonce in three children; this is a necessary probe, not full fidelity certification. Current protocol excludes failed restores from timing but counts them as failures; any amended failure or capture-amortization rule must be dated before collection. Resource cost and batch latency are different endpoints. A timing gain does not prove that process memory was needed: children must consume useful captured running state and the warm alternative must reconstruct equivalent usable state.}
\note{[1:10] The first experiment asks which execution mechanism we should choose. We compare restore after setup and a warm test with a cached template that already has dependencies. For three, six and twelve children, we interleave twenty batches per arm and measure the time until the last child produces its first test result. There are no model calls. We establish capture semantics first and expose failures and capture cost. A win informs the backend choice for this workload. To claim a benefit from live memory, we also need useful running state and an equivalent warm reconstruction baseline.}
\end{frame}
